% page 482 du fac-simile (image p-481 du scan)
% bloc 157.5 x 259.3 mm ; marge gauche 8.0 mm
\pgc{-12.700mm}{0.000mm}{
\l{\marge[80.760mm]{36.332mm}{\ornamento{12.928mm}{ornements/letroj/litero-Q.png}}}
\l{}
\l{}
\l{}
\l{}
\l{}
\l{}
\l{}
\l{}
\l{}
\l{}
\l{}
\l{\sou{qua}.\cel{2}(\sou{gram}.)\cel{2}I. (\sou{pronomo relativa}).\cel{2}Ne havas ipse senco ;}
\l{\cel{3}lokizita kom vorto rang-unesma di propoziciono qua plu-}
\l{\cel{3}precizigas antecedento tre nepreciza, e determinas lu :}
\l{\cel{3}(a) se \textquotedbl{}qua\textquotedbl{} ne preiresas nemediate da komo, la anteceden-}
\l{\cel{3}to esas la\cel{1}\sou{lasta}\cel{1}substantivo di la propoziciono preiranta ;}
\l{\cel{3}(b) se \textquotedbl{}qua\textquotedbl{} preiresas da komo, la antecedento esas la}
\l{\cel{3}\sou{pre-lasta}\cel{1}substantivo di la propoziciono preiranta. (Ta}
\l{\cel{3}regulo esas aplikita rigoroze en ica Radikaro). - II. (\sou{pro}-}
\l{\cel{3}\sou{nomo questionala}) : \textquotedbl{}Qua persono...?\textquotedbl{}.- III. (\sou{adjektivo}}
\l{\cel{3}\sou{questionala}) \textquotedbl{}Qua... ek la plura qui existas...?\textquotedbl{}}
\l{}
\l{\sou{quadro}.\cel{2}Petro-bloko qua esas taliita (o qua esos taliita)}
\l{\cel{3}por esor uzebla kom parto di domo, di edifico. - DI.}
\l{}
\l{\sou{quadrangula}. (\sou{geom.)}\cel{3}Qua havas quar anguli e quar lateri.}
\l{\cel{3}-\cel{1}\sou{Prismo, piramido quadriangula}\cel{1}: di qua quadrilatero for-}
\l{\cel{3}macas la bazo. -}
\l{}
\l{\sou{quadrant}o. (\sou{geom.})\cel{2}Quarimo di la cirklo-kurvo, t.e. 90.}
\l{\cel{3}- DEFIRS.}
\l{}
\l{\sou{quadrato}. (\sou{geom.}) Quadriangulo equilatera ed ortangula.}
\l{\cel{3}- DEFIRS.}
\l{}
\l{\sou{quadratika}.\cel{2}(\sou{mat.)}\cel{2}\sou{Equaciono quadratika}\cel{1}:\cel{2}equaciono alge-}
\l{\cel{3}brala di qua la grado egalesas 2-potenco e di qua la solvo}
\l{\cel{3}esas reduktebla ad olta di nombro finita de equacioni di}
\l{\cel{3}la grado duesma. -}
\l{}
\l{\sou{quadraturo}.\cel{2}(\sou{astron.)}\cel{3}Aspekto di du astri qui interdistas}
\l{\cel{3}per un quarimo di cirklokurvo. - DEFIRS.}
\l{}
\l{\sou{quadri...}\cel{2}(\sou{prefixo ciencala}) \textquotedbl{}Qua havas quar....-i\textquotedbl{}.}
\l{}
\l{\sou{quadrigo}. (en\cel{1}\sou{la epoki antiqua}). Charo adqua onu jungis quar}
\l{\cel{3}kavali sama-fronte. - DEFIS.}
\l{}
\l{\sou{quadriko}. (\sou{geom.}) Areo di la grado duesma, t.e. reprezentata}
\l{\cel{3}da equaciono di la grado duesma. - DEFIS.}
\l{}
\l{\sou{quadrilo}. I. (a) Trupo de kavalieri, dividita (ordinare) aden}
\l{\cel{3}quar grupi, qua figurantesas en karuselo ; (b) singla de}
\l{\cel{3}ta grupi, distingita per kolori, per kostumi interdifer-}
\l{\cel{3}anta. - II. (a) Nombro para de dui de danseruli e de dans-}
\l{\cel{3}erini qui figurantesas, ici regardante iti facie, en kon-}
\l{\cel{3}tradanso ; (b) la danso ipsa. - III. (\sou{teknol.})\cel{2}To quo}
\l{\cel{3}reprezentas quadrato, rombo, ajuroza o sen-ajura, en la}
\l{\cel{3}desegnuri di stofo, gipuro, pasmento, ec. DEFIRS.}
}
